\heading{高等数学A（II）-25春-期末}
\noteinfo{题目由用户提供，解答按本仓库现有版式整理。}

\begin{question}[正项级数的敛散性]

判别下列正项级数的敛散性：
\begin{enumerate}[label=(\arabic*),leftmargin=2.6em,itemsep=0.25em,topsep=0.2em]
\item $\sum_{n=1}^{\infty}\left(\ln \dfrac{1}{n}-\ln \sin \dfrac{1}{n}\right)$；
\item $\sum_{n=1}^{\infty}\dfrac{n!\arctan\!\bigl(2+(-1)^n\bigr)}{n^n}$；
\item $\sum_{n=1}^{\infty}\left(\arctan\dfrac{1}{n}\right)^\alpha\dfrac{1}{\ln(1+n+n^2)}\ (\alpha>0)$；
\item $\sum_{n=1}^{\infty}\left(\cos\dfrac{1}{\sqrt n}\right)^{n^2}$。
\end{enumerate}

\begin{solution}
\begin{enumerate}[label=(\arabic*),leftmargin=2.6em,itemsep=0.25em,topsep=0.2em]
\item 记
\[
a_n=\ln\frac{1/n}{\sin(1/n)}=\ln\frac{x}{\sin x}\Big|_{x=1/n}.
\]
当 $x\to0$ 时，
\[
\sin x=x-\frac{x^3}{6}+O(x^5),
\qquad
\frac{x}{\sin x}=1+\frac{x^2}{6}+O(x^4),
\]
从而
\[
\ln\frac{x}{\sin x}=\frac{x^2}{6}+O(x^4).
\]
取 $x=1/n$，得
\[
a_n=\frac{1}{6n^2}+O\!\left(\frac{1}{n^4}\right).
\]
故原级数与 $\sum 1/n^2$ 同敛散，所以
\ansmath{\sum_{n=1}^{\infty}\left(\ln \frac{1}{n}-\ln \sin \frac{1}{n}\right)\text{ 收敛。}}

\item 由于
\[
\arctan 1\le \arctan\!\bigl(2+(-1)^n\bigr)\le \arctan 3,
\]
故它只差一个正常数因子。设
\[
b_n=\frac{n!\arctan\!\bigl(2+(-1)^n\bigr)}{n^n}.
\]
则
\[
\frac{b_{n+1}}{b_n}
=\frac{\arctan(2+(-1)^{n+1})}{\arctan(2+(-1)^n)}
\left(\frac{n}{n+1}\right)^n.
\]
而
\[
\left(\frac{n}{n+1}\right)^n\to e^{-1},
\]
前面的比值有界，因此 $\limsup\limits_{n\to\infty} b_{n+1}/b_n<1$。由比值判别法知原级数收敛，即
\ansmath{\sum_{n=1}^{\infty}\frac{n!\arctan\!\bigl(2+(-1)^n\bigr)}{n^n}\text{ 收敛。}}

\item 设
\[
c_n=\left(\arctan\frac1n\right)^\alpha\frac{1}{\ln(1+n+n^2)}.
\]
由
\[
\arctan\frac1n\sim \frac1n,
\qquad
\ln(1+n+n^2)\sim 2\ln n
\]
可得
\[
c_n\sim \frac{1}{2n^\alpha\ln n}.
\]
故原级数与
\[
\sum_{n=2}^{\infty}\frac{1}{n^\alpha\ln n}
\]
同敛散，从而
\ansmath{\alpha>1\text{ 时收敛，}\ 0<\alpha\le1\text{ 时发散。}}

\item 用根值判别法。设
\[
d_n=\left(\cos\frac{1}{\sqrt n}\right)^{n^2},
\qquad
\sqrt[n]{d_n}=\left(\cos\frac{1}{\sqrt n}\right)^n.
\]
当 $x\to0$ 时，
\[
\ln(\cos x)=-\frac{x^2}{2}+o(x^2),
\]
于是
\[
n\ln\!\left(\cos\frac{1}{\sqrt n}\right)\to -\frac12,
\qquad
\left(\cos\frac{1}{\sqrt n}\right)^n\to e^{-1/2}<1.
\]
故原级数收敛，即
\ansmath{\sum_{n=1}^{\infty}\left(\cos\frac{1}{\sqrt n}\right)^{n^2}\text{ 收敛。}}
\end{enumerate}
\end{solution}
\end{question}

\begin{question}[含参数级数的敛散性]

根据 $\alpha$ 的取值讨论下列级数的敛散性；若是非正项级数，需进一步讨论条件收敛与绝对收敛：
\begin{enumerate}[label=(\arabic*),leftmargin=2.6em,itemsep=0.25em,topsep=0.2em]
\item $\sum_{n=2}^{\infty}\ln\!\left(1+\dfrac{(-1)^n}{n^\alpha}\right)$，$\alpha\in\mathbb R$；
\item $\sum_{n=1}^{\infty}\left(\sqrt[n]{\alpha}-\sqrt{1+\dfrac{1}{n}}\right)$，$\alpha>0$。
\end{enumerate}

\begin{solution}
\begin{enumerate}[label=(\arabic*),leftmargin=2.6em,itemsep=0.25em,topsep=0.2em]
\item 当 $\alpha\le0$ 时，奇数项中
\[
1+\frac{(-1)^n}{n^\alpha}=1-\frac{1}{n^\alpha}
\]
从某项起不为正，因此该级数本身无意义。下面只讨论 $\alpha>0$。

令
\[
u_n=\frac{(-1)^n}{n^\alpha}.
\]
由 Taylor 展开，
\[
\ln(1+u_n)=u_n-\frac{u_n^2}{2}+O(u_n^3)
=\frac{(-1)^n}{n^\alpha}-\frac{1}{2n^{2\alpha}}+O\!\left(\frac{1}{n^{3\alpha}}\right).
\]
于是原级数可分解为
\[
\sum \frac{(-1)^n}{n^\alpha}-\frac12\sum \frac{1}{n^{2\alpha}}
+\sum O\!\left(\frac{1}{n^{3\alpha}}\right).
\]
从而：
\begin{itemize}[leftmargin=2em,itemsep=0.2em,topsep=0.2em]
\item 当 $\alpha>1$ 时，$\sum |\ln(1+u_n)|$ 收敛，故绝对收敛；
\item 当 $\frac12<\alpha\le1$ 时，$\sum (-1)^n/n^\alpha$ 收敛，而 $\sum 1/n^{2\alpha}$ 也收敛，所以原级数收敛但不绝对收敛；
\item 当 $0<\alpha\le\frac12$ 时，$\sum 1/n^{2\alpha}$ 发散，且展开中的二次项为负，故原级数发散。
\end{itemize}
综上，
\ansmath{
\begin{cases}
\alpha>1,& \text{绝对收敛},\\[2mm]
\dfrac12<\alpha\le1,& \text{条件收敛},\\[2mm]
0<\alpha\le\dfrac12,& \text{发散},\\[2mm]
\alpha\le0,& \text{级数无意义}.
\end{cases}}

\item 记
\[
a_n=\alpha^{1/n}-\sqrt{1+\frac1n}.
\]
当 $n\to\infty$ 时，
\[
\alpha^{1/n}=e^{(\ln\alpha)/n}
=1+\frac{\ln\alpha}{n}+O\!\left(\frac{1}{n^2}\right),
\]
\[
\sqrt{1+\frac1n}=1+\frac{1}{2n}+O\!\left(\frac{1}{n^2}\right).
\]
故
\[
a_n=\frac{\ln\alpha-\frac12}{n}+O\!\left(\frac{1}{n^2}\right).
\]
于是：
\begin{itemize}[leftmargin=2em,itemsep=0.2em,topsep=0.2em]
\item 若 $\ln\alpha\ne \frac12$，即 $\alpha\ne \sqrt e$，则 $a_n\sim c/n$，级数发散；
\item 若 $\alpha=\sqrt e$，则首项消去，有 $a_n=O(1/n^2)$，故级数收敛。
\end{itemize}
因此
\ansmath{\sum_{n=1}^{\infty}\left(\sqrt[n]{\alpha}-\sqrt{1+\frac1n}\right)\text{ 当且仅当 }\alpha=\sqrt e\text{ 时收敛。}}
\end{enumerate}
\end{solution}
\end{question}

\begin{question}[幂级数展开与求和]

求函数
\[
f(x)=\arctan\frac{1-2x}{1+2x}
\]
在 $0$ 处的幂级数展开式，并求
\[
\sum_{n=0}^{\infty}\frac{(-1)^n}{2n+1}
\]
的和。

\begin{solution}
在 $|x|<\frac12$ 内，
\[
\arctan\frac{1-2x}{1+2x}=\frac{\pi}{4}-\arctan(2x),
\]
因为
\[
\tan\!\left(\frac{\pi}{4}-\arctan(2x)\right)=\frac{1-2x}{1+2x},
\]
且两边在 $x=0$ 时都等于 $\pi/4$。

又
\[
\arctan t=\sum_{n=0}^{\infty}\frac{(-1)^n t^{2n+1}}{2n+1},
\qquad |t|<1,
\]
取 $t=2x$，得
\[
f(x)=\frac{\pi}{4}-\sum_{n=0}^{\infty}\frac{(-1)^n(2x)^{2n+1}}{2n+1},
\qquad |x|<\frac12.
\]
即
\ansmath{
f(x)=\frac{\pi}{4}-\sum_{n=0}^{\infty}\frac{(-1)^n2^{2n+1}}{2n+1}x^{2n+1},
\qquad |x|<\frac12.}

再取 $x=\frac12$，则左边为
\[
f\!\left(\frac12\right)=\arctan 0=0,
\]
右边化为
\[
\frac{\pi}{4}-\sum_{n=0}^{\infty}\frac{(-1)^n}{2n+1}.
\]
故
\ansmath{\sum_{n=0}^{\infty}\frac{(-1)^n}{2n+1}=\frac{\pi}{4}.}
\end{solution}
\end{question}

\begin{question}[求和]

求级数
\[
\sum_{n=1}^{\infty}\frac{(-1)^n}{3n-1}.
\]

\begin{solution}
将其改写为
\[
\sum_{n=1}^{\infty}\frac{(-1)^n}{3n-1}
=-\sum_{n=0}^{\infty}\frac{(-1)^n}{3n+2}.
\]
利用
\[
\frac{1}{3n+2}=\int_0^1 x^{3n+1}\,dx,
\]
得
\[
\sum_{n=1}^{\infty}\frac{(-1)^n}{3n-1}
=-\int_0^1\sum_{n=0}^{\infty}(-1)^n x^{3n+1}\,dx
=-\int_0^1\frac{x}{1+x^3}\,dx.
\]
对
\[
\frac{x}{1+x^3}
\]
作部分分式分解：
\[
\frac{x}{1+x^3}
=-\frac{1}{3(x+1)}+\frac{x+1}{3(x^2-x+1)}.
\]
故
\[
\int \frac{x}{1+x^3}\,dx
=-\frac13\ln(x+1)+\frac16\ln(x^2-x+1)
+\frac{1}{\sqrt3}\arctan\frac{2x-1}{\sqrt3}+C.
\]
代入 $0$ 与 $1$，得
\[
\int_0^1\frac{x}{1+x^3}\,dx=-\frac13\ln2+\frac{\pi}{3\sqrt3}.
\]
因此
\ansmath{\sum_{n=1}^{\infty}\frac{(-1)^n}{3n-1}=\frac13\ln2-\frac{\pi}{3\sqrt3}.}
\end{solution}
\end{question}

\begin{question}[Fourier 级数与 $\sum 1/n^6$]

设 $f$ 是 $\mathbb R$ 上 $2\pi$ 周期的奇函数，在 $[0,\pi]$ 上
\[
f(x)=x(\pi-x).
\]
\begin{enumerate}[label=(\arabic*),leftmargin=2.6em,itemsep=0.25em,topsep=0.2em]
\item 证明：对任意 $x\in\mathbb R$，
\[
f(x)=\frac{8}{\pi}\sum_{n=1}^{\infty}\frac{\sin\bigl((2n-1)x\bigr)}{(2n-1)^3}.
\]
\item 求 $\sum_{n=1}^{\infty}\dfrac{1}{n^6}$。
\end{enumerate}

\begin{solution}
\begin{enumerate}[label=(\arabic*),leftmargin=2.6em,itemsep=0.25em,topsep=0.2em]
\item 由于 $f$ 为奇函数，其 Fourier 级数只有正弦项：
\[
f(x)\sim \sum_{n=1}^{\infty}b_n\sin nx,
\qquad
b_n=\frac{2}{\pi}\int_0^{\pi}x(\pi-x)\sin nx\,dx.
\]
分部积分两次可得
\[
b_n=\frac{4\bigl(1-(-1)^n\bigr)}{\pi n^3}.
\]
故偶数项系数为 $0$，奇数项系数为
\[
b_{2n-1}=\frac{8}{\pi(2n-1)^3}.
\]
于是
\ansmath{f(x)=\frac{8}{\pi}\sum_{n=1}^{\infty}\frac{\sin\bigl((2n-1)x\bigr)}{(2n-1)^3}.}

\item 由 Parseval 等式，
\[
\frac{1}{\pi}\int_{-\pi}^{\pi}|f(x)|^2\,dx=\sum_{n=1}^{\infty}b_n^2.
\]
因为 $f$ 是奇函数，
\[
\frac{1}{\pi}\int_{-\pi}^{\pi}|f(x)|^2\,dx
=\frac{2}{\pi}\int_0^{\pi}x^2(\pi-x)^2\,dx
=\frac{2}{\pi}\cdot \frac{\pi^5}{30}
=\frac{\pi^4}{15}.
\]
另一方面，
\[
b_{2n-1}^2=\frac{64}{\pi^2(2n-1)^6},
\qquad b_{2n}=0,
\]
故
\[
\frac{64}{\pi^2}\sum_{n=1}^{\infty}\frac{1}{(2n-1)^6}=\frac{\pi^4}{15},
\]
即
\[
\sum_{n=1}^{\infty}\frac{1}{(2n-1)^6}=\frac{\pi^6}{960}.
\]
设
\[
S=\sum_{n=1}^{\infty}\frac{1}{n^6},
\]
则
\[
\sum_{n=1}^{\infty}\frac{1}{(2n-1)^6}=S-\frac{1}{2^6}S=\frac{63}{64}S.
\]
因此
\[
\frac{63}{64}S=\frac{\pi^6}{960},
\]
从而
\ansmath{\sum_{n=1}^{\infty}\frac{1}{n^6}=\frac{\pi^6}{945}.}
\end{enumerate}
\end{solution}
\end{question}

\begin{question}[反常积分的收敛条件]

设
\[
\int_0^{\infty}\frac{x^n(1-e^{-x})}{(1+x)^m}\,dx
\]
收敛，求 $m,n$ 的取值范围。

\begin{solution}
分别考察 $x\to0$ 与 $x\to\infty$。

当 $x\to0$ 时，
\[
1-e^{-x}\sim x,
\qquad
(1+x)^m\sim 1,
\]
故被积函数满足
\[
\frac{x^n(1-e^{-x})}{(1+x)^m}\sim x^{n+1}.
\]
于是积分在 $0$ 附近收敛当且仅当
\[
n+1>-1,
\qquad\text{即}\qquad n>-2.
\]

当 $x\to\infty$ 时，
\[
1-e^{-x}\to1,
\qquad
(1+x)^m\sim x^m,
\]
故被积函数满足
\[
\frac{x^n(1-e^{-x})}{(1+x)^m}\sim x^{n-m}.
\]
于是积分在无穷远处收敛当且仅当
\[
n-m<-1,
\qquad\text{即}\qquad m-n>1.
\]

综上，
\ansmath{n>-2,\qquad m-n>1.}
\end{solution}
\end{question}

\begin{question}[广义积分的敛散性]

判断广义积分
\[
\int_1^{\infty}\frac{(\arctan x)\ln x}{x^2}\,dx
\]
的敛散性。

\begin{solution}
当 $x\ge1$ 时，
\[
0<\arctan x<\frac{\pi}{2},
\]
故
\[
0\le \frac{(\arctan x)\ln x}{x^2}
\le \frac{\pi}{2}\cdot \frac{\ln x}{x^2}.
\]
而
\[
\int_1^{\infty}\frac{\ln x}{x^2}\,dx
\]
收敛，所以由比较判别法可知原广义积分收敛。故
\ansmath{\int_1^{\infty}\frac{(\arctan x)\ln x}{x^2}\,dx\text{ 收敛。}}
\end{solution}
\end{question}

\begin{question}[恒等式证明与高斯积分]

设
\[
f(t)=\left(\int_0^t e^{-x^2}\,dx\right)^2,
\qquad
g(t)=\int_0^1\frac{e^{-t^2(1+x^2)}}{1+x^2}\,dx.
\]
证明 $f(t)+g(t)=\dfrac{\pi}{4}$，并由此计算
\[
\int_0^{\infty}e^{-x^2}\,dx.
\]

\begin{solution}
先求导。由链式法则，
\[
f'(t)=2e^{-t^2}\int_0^t e^{-x^2}\,dx.
\]
对 $g(t)$ 在积分号下求导，得
\[
g'(t)=\int_0^1\frac{-2t(1+x^2)e^{-t^2(1+x^2)}}{1+x^2}\,dx
=-2t e^{-t^2}\int_0^1 e^{-t^2x^2}\,dx.
\]
令 $u=tx$，则
\[
\int_0^1 e^{-t^2x^2}\,dx=\frac1t\int_0^t e^{-u^2}\,du.
\]
故
\[
g'(t)=-2e^{-t^2}\int_0^t e^{-u^2}\,du=-f'(t).
\]
因此
\[
\bigl(f(t)+g(t)\bigr)'=0,
\]
即 $f(t)+g(t)$ 为常数。

取 $t=0$，有
\[
f(0)=0,
\qquad
g(0)=\int_0^1\frac{dx}{1+x^2}=\frac{\pi}{4}.
\]
所以
\ansmath{f(t)+g(t)=\frac{\pi}{4}.}

再令 $t\to\infty$，由夹逼或控制收敛定理知 $g(t)\to0$，于是
\[
\left(\int_0^{\infty}e^{-x^2}\,dx\right)^2=\lim_{t\to\infty}f(t)=\frac{\pi}{4}.
\]
由于该积分为正，故
\ansmath{\int_0^{\infty}e^{-x^2}\,dx=\frac{\sqrt\pi}{2}.}
\end{solution}
\end{question}

